(1) 令 \(p_e(t)\) 为激发态概率，泵浦使基态以速率 \(P\) 跃迁到激发态，自发辐射以 \(\gamma\) 使激发态回到基态。因此 \(\dot p_e=P(1-p_e)-\gamma p_e=P-(P+\gamma)p_e\)，并保持 \(0\le p_e\le1\)。
(2) 稳态 \(p_e^{\rm ss}=P/(P+\gamma)\)，光子发射率 \(R=\gamma p_e^{\rm ss}=\gamma P/(P+\gamma)\)。弱泵浦 \(P\ll\gamma\) 时 \(R\simeq P\)；强泵浦 \(P\gg\gamma\) 时 \(R\to\gamma\)，发射率饱和。
(3) 光子探测后条件初态为基态 \(p_e(0|g)=0\)，解速率方程得 \(p_e(\tau|g)=P[1-e^{-(P+\gamma)\tau}]/(P+\gamma)\)。稳态二阶相关为 \(g^{(2)}(\tau)=p_e(|\tau||g)/p_e^{\rm ss}=1-e^{-(P+\gamma)|\tau|}\)，忽略探测器死时间。
(4) \(g^{(2)}(0)=0\)，因为一个二能级发射体无法同时放出两个自发辐射光子；\(|\tau|\to\infty\) 时相关趋一。反聚束恢复时间为 \((P+\gamma)^{-1}\)，随泵浦增强而缩短，表明非经典单光子发射的时间排斥。
